% 译自 arXiv:2306.13589v2 原 tex 第 882–1290 行
\subsection{同伦范畴}\label{ch2.4}


我们接着引入命题 \ref{quasilocalizing} 所需要的全部工具，这些工具最终将证实命题 \ref{derivedstructure} 所给出的关于 $\mathsf{D(A)}$ 的显式描述。

\begin{Def}
	设 $\mathsf{A}$ 为阿贝尔范畴。则其\textbf{同伦范畴}（homotopy category）\index{homotopy category} $\mathsf{H(A)}$ 由 $\text{obj}(\mathsf{H(A)})\coloneqq\text{obj}(\mathsf{Kom(A)})$ 给出，并把同伦的态射 $f^\bullet\sim g^\bullet\in\text{Hom}_{\mathsf{Kom(A)}}(X^\bullet, Y^\bullet)$ 与同一个 $f^\bullet=g^\bullet\in\text{Hom}_{\mathsf{H(A)}}(X^\bullet, Y^\bullet)$ 等同起来。
	
	全子范畴 $\mathsf{H^+(A)}, \mathsf{H^-(A)}, \mathsf{H^b(A)}\subset\mathsf{H(A)}$ 由相应的有界复形范畴类似地构造。
\end{Def}

与引理 \ref{komabelian} 中同样的论证表明，$\mathsf{H(A)}$ 显然是加性范畴（additive category）；这实际上已足以把上同调函子限制为定义 \ref{cohomfunc} 中所描述的函子 $H^n:\mathsf{H(A)}\rightarrow\mathsf{A}$。\footnote{注意，同伦的态射 $f^\bullet\sim g^\bullet$ 在 $\mathsf{A}$ 中给出相同的上同调态射，这与 $\mathsf{H(A)}$ 中 $f^\bullet=g^\bullet$ 的等同更加相容。} 此外，把链映射映为其模同伦的等价类的函子 $\mathsf{Kom(A)}\rightarrow\mathsf{H(A)}$ 是加性函子。而定义 \ref{quasiso} 中关于拟同构（quasi-isomorphism）的定义原封不动地搬到 $\mathsf{H(A)}$ 上依然成立。

% [Tutor Note: 易错预警] 这里「把同伦的态射等同起来」只是商掉链同伦，对象类完全不变，故 $\text{obj}(\mathsf{H(A)})=\text{obj}(\mathsf{Kom(A)})$；$H^n$ 之所以能下降到 $\mathsf{H(A)}$，依据的正是脚注中「同伦态射诱导相同的上同调态射」这一事实。切勿把 $\mathsf{H(A)}$ 与最终用拟同构局部化得到的 $\mathsf{D(A)}$ 混为一谈：在 $\mathsf{H(A)}$ 中拟同构一般并不可逆。

遗憾的是，正如 Verdier 在 \cite[Proposition II.1.3.6]{[Ver96]} 中所证明的，$\mathsf{H(A)}$ 仅在非常特殊的情形下才是阿贝尔范畴：\footnote{关于 $\mathsf{H(A)}$ 非阿贝尔的一个直接反例见 \url{https://math.stackexchange.com/questions/1280199/the-homotopy-category-of-complexes/1284823\#1284823}。}

\begin{Lem}\label{whenD(A)abelian}
	设 $\mathsf{A}$ 为阿贝尔范畴。则其同伦范畴 $\mathsf{H(A)}$ 是阿贝尔范畴当且仅当 $\mathsf{A}$ 是半单（semisimple）的\textup（参看注 \ref{semisimple}\textup）。
\end{Lem}

% [Tutor Note: 含糊点拨] 此处的「半单」按注 \ref{semisimple} 的含义，即任意短正合序列都分裂。该等价表明：只要 $\mathsf{A}$ 不是半单的，$\mathsf{H(A)}$ 中就会出现不可裂的短正合序列，因而「序列正合」这一概念不再够用——这正是下文改用三角形的动机。证明可参看作者所引 Verdier \cite[Proposition II.1.3.6]{[Ver96]} 及脚注中的反例链接。

因此，人们所期望的序列正合性概念在这里已经失效，应当由第 \ref{ch2.3} 节的三角形取而代之。而取道同伦范畴的好处则另有一层：

\begin{Thm}
	设 $\mathsf{A}$ 为阿贝尔范畴，并考虑类 $S\coloneqq\{\text{quasi-isomorphisms in }\mathsf{H(A)}\}$。则局部化（localization）$\mathsf{H(A)}[S^{-1}]$ \textup(按命题 \ref{derivedexists} 的证明中所构造的方式\textup) 典范同构于 $\mathsf{D(A)}$。类似地，对 $\#=+,-,\mathsf{b}$，$\mathsf{H^\#\!(A)}[S^{-1}]$ 典范同构于 $\mathsf{D^\#\!(A)}$。
\end{Thm}

\begin{proof}
	由于复合函子 $\mathsf{F}:\mathsf{Kom(A)}\rightarrow\mathsf{H(A)}\rightarrow\mathsf{H(A)}[S^{-1}]$ 依构造把拟同构映为同构，故由命题 \ref{derivedexists} 存在唯一函子 $\mathsf{G}:\mathsf{D(A)}\rightarrow\mathsf{H(A)}[S^{-1}]$ 使 $\mathsf{F}=\mathsf{G}\circ\mathsf{L}$。对任意 $X\in\text{obj}(\mathsf{Kom(A)})=\text{obj}(\mathsf{D(A)})$ 有 $\mathsf{G}(X)=(\mathsf{G}\circ\mathsf{L})(X)=\mathsf{F}(X)=X\in\text{obj}(\mathsf{H(A)}[S^{-1}])$，因此 $\mathsf{G}$ 在对象上是双射。此外，$\mathsf{H(A)}[S^{-1}]$ 中的任一态射类 $[f^\bullet]$ 都在 $\mathsf{H(A)}$ 中有代表元 $f^\bullet$，从而按同伦范畴的定义必在 $\mathsf{Kom(A)}$ 中也有代表元，于是 $[f^\bullet]=\mathsf{F}(f^\bullet)=(\mathsf{G}\circ\mathsf{L})(f^\bullet)=\mathsf{G}([\tilde{f}^\bullet])$，这证明了 $\mathsf{G}$ 在态射上是满的。
	
	关于态射上的单射性，注意 $f^\bullet\sim g^\bullet\in\text{Hom}_{\mathsf{Kom(A)}}(X^\bullet, Y^\bullet)$ 蕴涵 $\mathsf{L}(f^\bullet)=\mathsf{L}(g^\bullet)\in\text{Hom}_{\mathsf{D(A)}}(X^\bullet, Y^\bullet)$。我们把这一断言明确证明出来。依假设，对某个同伦 $k^\bullet:X^\bullet\rightarrow Y^{\bullet-1}$ 有 $f^n=g^n+d_Y^{n-1}\circ k^n + k^{n+1}\circ d_X^n$。由此得到态射 $c(k)^\bullet:C(f)^\bullet=X^{\bullet+1}\oplus Y^\bullet\rightarrow C(g)^\bullet=X^{\bullet+1}\oplus Y^\bullet$，$c(k)^n(x^{n+1},y^n)\coloneqq(x^{n+1},y^n+k^{n+1}(x^{n+1}))$，它实际上是链映射，因为 
	\begin{align*}
		\big(d_{C(g)}^n\!\circ c(k)^n)&(x^{n+1},y^n\big)\!=\!\big(\!-d_X^{n+1}(x^{n+1}),g^{n+1}(x^{n+1})+d_Y^n(y^n+k^{n+1}(x^{n+1}))\big)\\
		&=\big(-d_X^{n+1}(x^{n+1}), f^{n+1}(x^{n+1})+d_Y^n(y^n)+k^{n+2}(-d_X^{n+1}(x^{n+1}))\big)\\
		&= c(k)^{n+1}\big(-d_X^{n+1}(x^{n+1}),f^{n+1}(x^{n+1})+d_Y^n(y^n)\big)\\
		&=\big(c(k)^{n+1}\circ d_{C(f)}^n\big)(x^{n+1},y^n).
	\end{align*}
	类似地，$k^\bullet$ 定义链映射 $z(k)^\bullet:Z(f)^\bullet\rightarrow Z(g)^\bullet$，$z(k)^n(x^n,x^{n+1},y^n)\coloneqq (x^n,x^{n+1},y^n+k^{n+1}(x^{n+1}))$。
	
	$f^\bullet$ 与 $g^\bullet$ 的正合行（如同 \eqref{distinguidiag} 中最上面的一行）合起来给出图
	\begin{equation*}
		\begin{tikzcd}
			0\arrow[r] & Y^\bullet\arrow[r,"i_2^\bullet"]\arrow[d,"\text{id}_Y^\bullet"'] & C(f)^\bullet\arrow[r,"p_1^\bullet"]\arrow[d,"c(k)^\bullet"] & X[1]^\bullet\arrow[r]\arrow[d,"\text{id}_{X[1]}^\bullet"] & 0 \\
			0\arrow[r] & Y^\bullet\arrow[r,"i_2^\bullet"'] & C(g)^\bullet\arrow[r,"p_1^\bullet"'] & X[1]^\bullet\arrow[r] & 0
		\end{tikzcd}\quad,
	\end{equation*}
	依 $c(k)^\bullet$ 的定义，它是交换的。由于这两行在 $\mathsf{Kom(A)}$ 中都是短正合的，这在上同调中诱导出一个图，其中所有竖直箭头仍是恒等态射，唯一的例外是 $H^n(c(k)^\bullet):H^n(C(f)^\bullet)\rightarrow H^n(C(g)^\bullet)$，而它由五引理（Five-Lemma）知为同构。因此，$c(k)^\bullet$ 是拟同构。这一点又有助于证明 $z(k)^\bullet$ 也是拟同构，后者可类似地由引理 \ref{p3i3quiso} 的中间行推出。
	
	最后，考虑图
	\[
	\begin{tikzcd}
		& X^\bullet\arrow[r, "\text{id}_X^\bullet"]\arrow[dl, "f^\bullet"']\arrow[d, "i_1^\bullet"] & X^\bullet\arrow[d, "i_1^\bullet"']\arrow[dr, "g^\bullet"] & \\
		Y^\bullet\arrow[r, "i_3^\bullet"'] & Z(f)^\bullet\arrow[r, "z(k)^\bullet"'] & Z(g)^\bullet\arrow[r, "p_3^\bullet"'] & Y^\bullet
	\end{tikzcd}\quad,
	\]
	它只在方块与右三角形处交换（其中 $p_3^\bullet$ 是通常的「修正」投影）。左三角形在 $\mathsf{Kom(A)}$ 中不交换\big(显然 $i_3^n(f^n(x^n))=(0,0,f^n(x^n))\neq(x^n,0,0)=i_1^n(x^n)$\big)，但在 $\mathsf{D(A)}$ 中它\textit{确实}交换，因为 $f^\bullet=p_3^\bullet\circ i_1^\bullet$（见右三角形）在 $\mathsf{L}:\mathsf{Kom(A)}\rightarrow\mathsf{D(A)}$ 下变成 $\mathsf{L}(f^\bullet)=\mathsf{L}(p_3^\bullet)\circ\mathsf{L}(i_1^\bullet)$，从而 $\mathsf{L}(i_1^\bullet)=\mathsf{L}(p_3^\bullet)^{-1}\circ\mathsf{L}(f^\bullet)=\mathsf{L}(i_3^\bullet)\circ\mathsf{L}(f^\bullet)$（再次用到引理 \ref{p3i3quiso}）。借助该图的完全交换性以及 $p_3^\bullet\circ z(k)^\bullet\circ i_3^\bullet=\text{id}_Y^\bullet$，我们最终推出 $\mathsf{L}(f^\bullet)=\mathsf{L}(p_3^\bullet\circ z(k)^\bullet\circ i_3^\bullet\circ f^\bullet)=\mathsf{L}(g^\bullet)$。
	
	因此，$\mathsf{G}$ 在态射上是单射：假设在 $\mathsf{H(A)}[S^{-1}]$ 中 $[f^\bullet]=[g^\bullet]$，亦即它们在最坏情形下于 $\mathsf{Kom(A)}$ 中有同伦的代表元 $f^\bullet\sim g^\bullet$，则我们刚刚证明了 $\mathsf{D(A)}$ 中 $\mathsf{L}(f^\bullet)=\mathsf{L}(g^\bullet)$，这正是所要的。我们断定 $\mathsf{G}$ 在对象与态射上都是双射，从而建立了 $\mathsf{H(A)}[S^{-1}]$ 与 $\mathsf{D(A)}$ 之间的等价，以及 $\mathsf{H^\#\!(A)}[S^{-1}]$ 与 $\mathsf{D^\#\!(A)}$ 之间的等价。
\end{proof}

% [Tutor Note: 含糊点拨] 「由五引理知为同构」一步省略了细节：所诱导的上同调图两行仍正合，左右两条竖直箭头为恒等，故中间箭头被夹逼成同构；随后对 $z(k)^\bullet$ 的拟同构性用的是引理 \ref{p3i3quiso} 中「$Z(f)^\bullet$ 的中间行」对应的相应交换图，作者仅以「类似地」带过。若需逐步验证，可自行把 \eqref{distinguidiag} 中关于 $C(f)^\bullet$ 的行替换为 $Z(f)^\bullet$ 的行后重复上述追图。

借助刚刚证明的结论，我们不妨干脆把阿贝尔范畴 $\mathsf{A}$ 的导出范畴\textit{定义}为 $\mathsf{D(A)}\coloneqq\mathsf{H(A)}[S^{-1}]$。现在，如以下结果所保证的，命题 \ref{derivedstructure} 可以应用了。

\begin{Pro}\label{quasilocalizing}
	设 $\mathsf{A}$ 为阿贝尔范畴。则态射类 $S\coloneqq\{\text{quasi-isomorphisms in }\mathsf{H(A)}\}$ 在 $\mathsf{H(A)}$ 中是局部化的\textup(类似地，$S^\#\coloneqq\{\text{quasi-isomorphisms in }\mathsf{H^\#\!(A)}\}$ 在 $\mathsf{H^\#\!(A)}$ 中是局部化的\textup)。 
\end{Pro}   

% [Tutor Note: 术语提醒] 「$S$ 在 $\mathsf{H(A)}$ 中是局部化类（localizing class）」是指定义 \ref{locclass} 的三条（含 Ore 条件型的两条补全性质）均成立，因此命题 \ref{derivedstructure} 关于分式演算的整套机器可原样使用，而无需额外验证饱和性。另请注意作者把 $\mathsf{H(A)}[S^{-1}]$ 就取作 $\mathsf{D(A)}$ 的\textit{定义}，故下文凡出现 $\mathsf{H(A)}[S^{-1}]$ 处均指同一个范畴。

\begin{proof}
	我们必须验证定义 \ref{locclass} 的三条均成立。a. 条是\break 显然的：对所有 $X^\bullet\in\text{obj}(\mathsf{H(A)})$ 有 $\text{id}_X^\bullet\in S$，且由每个 $H^n$ 的函子性，$s^\bullet,t^\bullet\in S$ 蕴涵 $s^\bullet\circ t^\bullet\in S$。
	
	关于 b. 条，任取 $\mathsf{H(A)}$ 中的三元组 $Z^\bullet\xrightarrow{s^\bullet}Y^\bullet\xleftarrow{f^\bullet}X^\bullet$，其中 $s^\bullet\in S$。目标是找到补全的三元组，使图 \eqref{localsquare} 的左边交换。设 $i_2^\bullet:Y^\bullet\rightarrow C(s)^\bullet=Z^{\bullet+1}\oplus Y^\bullet$ 为第二个嵌入，于是 $i_2^\bullet\circ f^\bullet:X^\bullet\rightarrow C(s)^\bullet$，\break 并取 $W^\bullet\coloneqq C(i_2\circ f)[-1]^\bullet= X^\bullet\oplus C(s)^{\bullet-1}=X^\bullet\oplus Z^\bullet\oplus Y^{\bullet-1}$ 与 $t^\bullet\coloneqq p_1^\bullet:W^\bullet\rightarrow X^\bullet$。它们纳入下图
	\begin{equation}\label{messydiag}
		\begin{tikzcd}
			{C(i_2\circ f)[-1]^\bullet}\arrow[r, dashed, "{p_1^\bullet}"]\arrow[d, dashed, "g^\bullet"'] & X^\bullet\arrow[r, "{i_2^\bullet\circ f^\bullet}"]\arrow[d, "f^\bullet"'] & C(s)^\bullet\arrow[r, "{i_2^\bullet}"]\arrow[d, "{\text{id}_{C(s)}^\bullet}"] & C(i_2\circ f)^\bullet\arrow[d, "{g[1]}^\bullet"] \\
			Z^\bullet\arrow[r, "s^\bullet"'] & Y^\bullet\arrow[r, "{i_2^\bullet}"'] & C(s)^\bullet\arrow[r, "{p_1^\bullet}"'] & Z[1]^\bullet
		\end{tikzcd}\quad,
	\end{equation}
	其中虚线箭头提示了所期望的态射。定义 $g^\bullet:W^\bullet\rightarrow Z^\bullet$ 为 $g^n(x^n,z^n,y^{n-1})$ $\coloneqq -z^n$，则左方块模同伦 $k^\bullet:W^\bullet\rightarrow Y^{\bullet-1}$（由 $k^n(x^n,z^n,y^{n-1})\coloneqq -y^{n-1}$ 给出）交换：事实上我们计算 $\big(k^{n+1}\circ d_W^n + d_Y^{n-1}\circ k^n\big)(x^n,z^n,y^{n-1})=-(-f^n(x^n)-s^n(z^n)-d_Y^{n-1}(y^{n-1}))+d_Y^{n-1}(-y^{n-1})=f^n(x^n) + s^n(z^n)=\big(f^n\circ p_1^n -s^n\circ g^n\big)(x^n,z^n,y^{n-1})$，其中第一个等号用了 $W^\bullet$ 的微分的定义，
	\begin{align*}
		d_W^n=-d_{C(i_2\circ f)}^{n-1}=\!-\!\begin{pmatrix} -d_X^n & 0 \\ i_2^n\circ f^n & d_{C(s)}^{n-1} \end{pmatrix} \!=\! \begin{pmatrix}
			d_X^n & 0 & 0 \\ 0 & d_Z^n & 0 \\ -f^n & -s^n & -d_Y^{n-1}
		\end{pmatrix}\!:W^n\rightarrow W^{n+1},
	\end{align*}
	从而 $d_W^n(x^n,z^n,y^{n-1})=\big(d_X^n(x^n), d_Z^n(z^n),-f^n(x^n)-s^n(z^n)-d_Y^{n-1}(y^{n-1})\big)$。因此，左方块在 $\mathsf{H(A)}$ 中交换。最后，我们论证 $p_1^\bullet$ 是拟同构。把注 \ref{C(f)acyclic} 应用于图 \eqref{messydiag} 的第二行可知 $C(s)^\bullet$ 无环。但第一行本身就是一个区分三角形（distinguished triangle），故定理 \ref{longexactcohomtriangles} 给出长正合序列
	\[
	...\rightarrow H^{n-1}(C(s)^\bullet)\cong 0 \rightarrow H^{n}(W^\bullet)\xrightarrow{H^n(p_1^\bullet)} H^{n}(X^\bullet)\rightarrow H^{n}(C(s)^\bullet)\cong 0\rightarrow...\,,
	\]
	这又告诉我们，对所有 $n\in\mathbb{Z}$，$H^n(p_1^\bullet)$ 必为同构，即 $p_1^\bullet\in S$。b. 条于是得到满足（对应于 \eqref{localsquare} 右边的第二个条件可类似地证明）。
	
	关于 c. 条，设 $f^\bullet\in\text{Hom}_\mathsf{H(A)}(X^\bullet,Y^\bullet)$，并假设存在 $S$ 中的 $s^\bullet:Y^\bullet\rightarrow Z^\bullet$ 使 $s^\bullet\circ f^\bullet=0$，意即 $s^\bullet\circ f^\bullet\sim 0$ 在 $\mathsf{Kom(A)}$ 中经由某个 $k^\bullet:X^\bullet\rightarrow Z^{\bullet-1}$ 而同伦于零（从而满足 $d_Z^{n-1}\circ k^n= s^n\circ f^n -k^{n+1}\circ d_X^n$）。于是我们要找 $S$ 中的某个 $t^\bullet:W^\bullet\rightarrow X^\bullet$ 使得 $f^\bullet\circ t^\bullet=0$。为此，考虑
	\begin{equation*}
		\begin{tikzcd}
			C(s)[-1]^\bullet\arrow[r, "{p_1^\bullet}"] & Y^\bullet\arrow[r, "s^\bullet"] & Z^\bullet\arrow[r, "{i_2^\bullet}"] & {C(s)^\bullet = Y^{\bullet+1}\oplus Z^\bullet} \\
			C(s)[-1]^\bullet\arrow[u, "{\text{id}_{C(s)[-1]}^\bullet}"] & X^\bullet\arrow[l, "g^\bullet"]\arrow[u, "f^\bullet"]\arrow[ur, "0"] & C(g)[-1]^\bullet\arrow[l, dashed, "t^\bullet"] & C(s)[-2]^\bullet=C(s)^{\bullet-2}\arrow[l, "i_2^\bullet"]\\
		\end{tikzcd}\quad,
		\vspace*{-0.3cm}
	\end{equation*}
	其中使方块交换的态射 $g^\bullet$ 由 $g^n(x^n)\coloneqq(f^n(x^n),-k^n(x^n))\in Y^n\oplus Z^{n-1}$ 给出，它是一个良定义的链映射： 
	\begin{align*}
		-d_{C(s)}^{n-1}&(g^n(x^n))=\big(d_Y^n(f^n(x^n)),-s^n(f^n(x^n))+d_Z^{n-1}(k^n(x^n))\big) \\
		&=\big(f^{n+1}(d_X^n(x^n)),-k^{n+1}(d_X^n(x^n))-d_Z^{n-1}(k^n(x^n))+d_Z^{n-1}(k^n(x^n))\big) \\
		&=g^{n+1}(d_X^n(x^n)).
	\end{align*}
	取 $W^\bullet\coloneqq C(g)[-1]^\bullet = X^\bullet\oplus C(s)^{\bullet-2}=X^\bullet\oplus Y^{\bullet-1}\oplus Z^{\bullet-2}$ 以及 $t^\bullet\coloneqq p_1^\bullet:W^\bullet\rightarrow X^\bullet$，$t^n(x^n,y^{n-1},z^{n-2})= x^n$。则如所愿 $f^\bullet\circ t^\bullet=p_1^\bullet\circ g^\bullet\circ t^\bullet=0$，因为底行正合。此外，由注 \ref{C(f)acyclic}，$C(s)^\bullet$ 无环，故底行中区分三角形在上同调中诱导的长正合序列蕴涵 $t^\bullet\in S$。反之，即由适当的 $t^\bullet$ 的存在性推出适当的 $s^\bullet$ 的存在性，可类似地证明。于是 c. 条成立，证毕。 
\end{proof} 


\subsection{导出范畴的结构}\label{ch2.5}


鉴于命题 \ref{quasilocalizing}，我们再把导出范畴的一般结构详细陈述一遍。

\begin{Def}\label{dercatdef}
	设 $\mathsf{A}$ 为阿贝尔范畴，其同伦范畴为 $\mathsf{H(A)}$，局部化类为 $S\coloneqq\{\text{quasi-isomorphisms in }\mathsf{H(A)}\}\subset\mathsf{H(A)}$。则其\textit{导出范畴} $\mathsf{D(A)}=\mathsf{H(A)}[S^{-1}]$ 具有如下结构：
	\begin{itemize}[leftmargin=0.5cm]
		\renewcommand{\labelitemi}{\textendash}
		\item $\text{obj}(\mathsf{D(A)})=\text{obj}(\mathsf{Kom(A)})$，因此对象是上链复形 $(X^\bullet,d_X^\bullet)$，其中每个 $X^n\in\text{obj}(\mathsf{A})$；
		
		\item 态射 $\varphi\in\text{Hom}_\mathsf{D(A)}(X^\bullet,Y^\bullet)$ 是屋顶图（roof）的等价类 $[(s^\bullet,f^\bullet)]$：
		\begin{equation}
			\begin{tikzcd}
				& {X'}^\bullet\arrow[dl, "s^\bullet"']\arrow[dr, "f^\bullet"] & \\
				X^\bullet & & Y^\bullet
			\end{tikzcd}\quad,
		\end{equation}
		其中 $f^\bullet\in\text{Hom}_{\mathsf{H(A)}}(X'^\bullet,Y^\bullet)$ 且 $s^\bullet\in S$。定义所用的等价关系由图 \eqref{equivroofs} 给出，态射的复合则按 \eqref{roofcompo} 进行。
	\end{itemize}
	导出范畴带有\textit{局部化函子}\footnote{与加性函子 $\mathsf{Kom(A)}\rightarrow\mathsf{H(A)}$ 复合，便得到命题 \ref{derivedexists} 中先前的局部化 $\mathsf{L}:\mathsf{Kom(A)}\rightarrow\mathsf{D(A)}$，后者随之传递其泛性质。} $\mathsf{L}:\mathsf{H(A)}\rightarrow\mathsf{D(A)}$，它在对象上是恒等，并把任意 $f^\bullet\in\text{Hom}_\mathsf{H(A)}(X^\bullet,Y^\bullet)$ 映为类 $[(\text{id}_X^\bullet,f^\bullet)]\in\text{Hom}_\mathsf{D(A)}(X^\bullet,Y^\bullet)$。此外它还满足：对任意 $s^\bullet\in S$，$\mathsf{L}(s^\bullet)$ 是同构；并且若任何其他函子 $\mathsf{F}:\mathsf{H(A)}\rightarrow\mathsf{D}$ 也如此作用，则它分解为 $\mathsf{F}=\mathsf{G}\circ\mathsf{L}$，其中 $\mathsf{G}:\mathsf{D(A)}\rightarrow\mathsf{D}$ 唯一。
	各变体 $\mathsf{D^\#\!(A)}$ 具有类似的结构。
\end{Def}

% [Tutor Note: 术语提醒] 「屋顶图」的约定是：准同构 $s^\bullet$ 一律放在左边、指向来源对象 $X^\bullet$，真正的态射 $f^\bullet$ 放在右边；$\mathsf{D(A)}$ 中的态射是这种图形的等价类，而非某个具体的链映射。请特别注意 $s^\bullet$ 的方向——若把箭头画反成 $X^\bullet\to X'^\bullet$，则局部化后复合的次序与 \eqref{roofcompo} 全部颠倒，这一约定是本节所有计算的基准。

由引理 \ref{whenD(A)abelian} 我们知道，$\mathsf{D(A)}$ 一般不是阿贝尔范畴。但至少：

\begin{Lem}\label{D(A)additive}
	设 $\mathsf{A}$ 为阿贝尔范畴。则其导出范畴 $\mathsf{D(A)}$ 是加性范畴。
\end{Lem}

\begin{proof}
	我们必须说明如何把由屋顶图 $X^\bullet\xleftarrow{s_k^\bullet}X_k'^\bullet\xrightarrow{f_k^\bullet}Y^\bullet$（$k=1,2$）表示的态射 $\varphi_k\in\text{Hom}_\mathsf{D(A)}(X^\bullet,Y^\bullet)$ 相加。为此，我们须找到一个「共同分母」，即让它们共用同一个拟同构。这可由下图实现：
	\begin{equation}
		\begin{tikzcd}
			& & X''^\bullet\arrow[dl, dashed, "r_1^\bullet"']\arrow[dr, dashed, "r_2^\bullet"]\arrow[dd, "t^\bullet"] & & \\
			& X_1'^\bullet\arrow[dl, "f_1^\bullet"']\arrow[dr, "s_1^\bullet"'] & & X_2'^\bullet\arrow[dl, "s_2^\bullet"]\arrow[dr, "f_2^\bullet"] & \\
			Y^\bullet & & X^\bullet & & Y^\bullet 
		\end{tikzcd}\qquad,
	\end{equation}
	其中交换方块是利用图 \eqref{localsquare} 构造的；这告诉我们 $r_1^\bullet\in S$，从而 $r_2^\bullet$ 也必为拟同构。于是 $t^\bullet\coloneqq s_1^\bullet\circ r_1^\bullet= s_2^\bullet\circ r_2^\bullet\in S$ 就是我们所求的共同分母；再令 $g_k^\bullet\coloneqq f_k^\bullet\circ r_k^\bullet$，便知 $\varphi_k$ 也可由屋顶图 $X^\bullet\xleftarrow{t^\bullet}X''^\bullet\xrightarrow{g_k^\bullet}Y^\bullet$ 表示。然后我们令 $(\varphi_1+\varphi_2)\in\text{Hom}_\mathsf{D(A)}(X^\bullet,Y^\bullet)$ 为由屋顶图 $X^\bullet\xleftarrow{t^\bullet}X''^\bullet\xrightarrow{g_1^\bullet+g_2^\bullet}Y^\bullet$ 所表示的类。仔细验证可知，它只依赖于 $\varphi_1,\varphi_2$，而与所选的代表元无关。
	
	上述构造的对称性显然蕴涵 $\varphi_1+\varphi_2=\varphi_2+\varphi_1$；我们还记得，恒等态射可平凡地由屋顶图 $X^\bullet\xleftarrow{\text{id}_X^\bullet}X^\bullet\xrightarrow{\text{id}_X^\bullet}X^\bullet$ 表示。于是我们认出了群 $\text{Hom}_\mathsf{D(A)}(X^\bullet,Y^\bullet)$ 的阿贝尔群结构。复合的双加性只需（大量地）正确作图即可看出。因此，$\mathsf{D(A)}$ 是预加性范畴（preadditive category）。
	
	$\mathsf{D(A)}$ 的零对象就是平凡复形，而事实上从它出发或指向它的任何态射都必为零。\footnote{顺便观察到，$[(s^\bullet,f^\bullet)]=0\in\text{Hom}_\mathsf{D(A)}(X^\bullet,Y^\bullet)$ 当且仅当存在 $S$ 中的某个 $s^\bullet:Y^\bullet\rightarrow Z^\bullet$（此时必非零）使 $s^\bullet\circ f^\bullet=0$，或等价地，存在 $S$ 中的某个 $t^\bullet:W^\bullet\rightarrow X'^\bullet$ 使 $f^\bullet\circ t^\bullet=0$。一个有效的代表元是 $X^\bullet\xleftarrow[\text{id}_X^\bullet]{}X^\bullet\xrightarrow[0]{}Y^\bullet$。} 最后，关于公理 (A3)，给定两个对象 $X_1^\bullet,X_2^\bullet\in\text{obj}(\mathsf{D(A)})$ 及其直和 $Y^\bullet=X_1^\bullet\oplus X_2^\bullet\in\text{obj}(\mathsf{D(A)})$（以显然的方式定义），我们令 $i_k^\bullet\in\text{Hom}_\mathsf{H(A)}(X_k^\bullet,Y^\bullet)$ 由屋顶图 $X_k^\bullet\xleftarrow{\text{id}_{X_k}^\bullet} X_k^\bullet\xrightarrow{i_k^\bullet}Y^\bullet$ 表示，令 $p_k^\bullet\in\text{Hom}_\mathsf{H(A)}(Y^\bullet,X_k^\bullet)$ 由 $Y^\bullet\xleftarrow{\text{id}_Y^\bullet} Y^\bullet\xrightarrow{p_k^\bullet}X_k^\bullet$ 表示。于是 \eqref{directprodsum} 所描述的必要关系容易验证；例如
	\[
	\begin{tikzcd}[column sep=0.5cm]
		& & X_1^\bullet\arrow[dl, "\text{id}_{X_1}^\bullet"']\arrow[dr, "i_1^\bullet"] & & \\
		& X_1^\bullet\arrow[dl, "\text{id}_{X_1}^\bullet"']\arrow[dr, "i_1^\bullet"'] & & Y^\bullet\arrow[dl, "\text{id}_Y^\bullet"]\arrow[dr, "p_k^\bullet"] & \\
		X_1^\bullet\arrow[rr, squiggly, "i_1^\bullet"'] & & Y^\bullet\arrow[rr, squiggly, "p_k^\bullet"'] & & X_k^\bullet 
	\end{tikzcd}\;\;\;\;\;\;
	\begin{tikzcd}[column sep=0.5cm]
		& & Y^\bullet\arrow[dl, "\text{id}_Y^\bullet"']\arrow[dr, "\text{id}_Y^\bullet"]\arrow[dd, "\text{id}_Y^\bullet"] & & \\
		& Y^\bullet\arrow[dl, "i_1^\bullet\circ p_1^\bullet"']\arrow[dr, "\text{id}_Y^\bullet"'] & & Y^\bullet\arrow[dl, "\text{id}_Y^\bullet"]\arrow[dr, "i_2^\bullet\circ p_2^\bullet"] & \\
		Y^\bullet & & Y^\bullet\arrow[ll, squiggly, "i_1^\bullet\circ p_1^\bullet"]\arrow[rr, squiggly, "i_2^\bullet\circ p_2^\bullet"'] & & Y^\bullet 
	\end{tikzcd}	
	\]
	告诉我们 $p_1^\bullet\circ i_1^\bullet=\text{id}_{X_1}^\bullet$、$p_2^\bullet\circ i_1^\bullet=0$ 以及 $i_1^\bullet\circ p_1^\bullet + i_2^\bullet\circ p_2^\bullet = \text{id}_Y^\bullet$ 在代表元层面也成立，从而在 $\mathsf{D(A)}$ 中成立。因此，导出范畴是加性范畴。
\end{proof}

% [Tutor Note: 记号提醒] 此处作者把 $i_k^\bullet,p_k^\bullet$ 写成 $\text{Hom}_\mathsf{H(A)}$ 中的态射，但它们是由屋顶图（即分式演算的像）给出的，严格说来应视为 $\text{Hom}_\mathsf{D(A)}$ 中的态射；按上下文应读作后者的像。另注意第一幅图中打波浪线的箭头标签写作 $p_k^\bullet$（一般指标 $k$），下文的验证分别取 $k=1$ 与 $k=2$ 两条关系，故不宜把它误读为某个固定的 $p_1^\bullet$。

最后我们说明如何把阿贝尔范畴看作其导出范畴的全子范畴。

\begin{Def}
	设 $\mathsf{A}$ 为阿贝尔范畴。对每个 $k\in\mathbb{Z}$，任意 $X\in\text{obj}(\mathsf{A})$ 都定义 $\mathsf{Kom^\#\!(A)}$（相应地 $\mathsf{H^\#\!(A)}$ 或 $\mathsf{D^\#\!(A)}$，其中 $\#=\emptyset,+,-,\mathsf{b}$）中形如\footnote{从上下文应当能看清 $X[n]^\bullet$ 指的是 $n$ 次（可能非平凡）的平移复形，还是由一个对象诱导的复形。}
	\[
	X[-k]^\bullet\coloneqq\big(...\rightarrow 0\rightarrow 0\rightarrow X\rightarrow 0\rightarrow 0\rightarrow ...\big)\,.
	\]
	的复形，其中 $X$ 位于第 $k$ 个位置。我们称它为 $X$ 诱导的\textbf{$k$-复形}\index{kcomplex@$k$-complex}，它除 $H^k(X[-k]^\bullet)\cong X$ 外无环。在注 \ref{0-complex} 中我们已经遇到过 $0$-复形 $X[0]^\bullet$。
	
	反之，若 $X^\bullet\in\mathsf{Kom^\#\!(A)}$（或属于其他复形范畴）在第 $k$ 个位置之外无环，即对所有 $j\neq k$ 有 $H^j(X^\bullet)=0$（人们也说 $X^\bullet$ \textit{集中在上同调次数 $k$}），则称它为\textbf{$H^k$-复形}\index{hk@$H^k$-complex}。因此，$k$-复形都是 $H^k$-复形。 
\end{Def}

% [Tutor Note: 易错预警] 记号 $X[-k]^\bullet$ 极易看错：$X$ 位于第 $k$ 个位置，而该复形记作 $X[-k]^\bullet$，即平移的编号与位置的编号符号相反；这是为了与 $\mathsf{D(A)}$ 中 $X[k]=X[0][k]$ 的平移约定相容。此外请与 $\mathsf{J}(X)=X[0]^\bullet$ 相区分：后者才是对象 $X$ 本身作为复形时的像，而 $X[-k]^\bullet$ 是它的 $k$ 次平移；脚注提醒的「$X[n]^\bullet$ 究竟指平移复形还是对象诱导复形」的歧义正源于此。

\begin{Pro}\label{AinD(A)}
	设 $\mathsf{A}$ 为阿贝尔范畴。则函子 $\mathsf{J'}\coloneqq \mathsf{L\circ J}:\mathsf{A}\rightarrow\mathsf{D^\#\!(A)}$ 是全忠实（fully faithful）的\textup(再次参看注 \ref{0-complex}\textup)。特别地，$\mathsf{A}$ 等价于\footnote{即 $\mathsf{J'}:\mathsf{A}\rightarrow\mathsf{D_0^\#\!(A)}$ 还是本质满的，意思是任意 $H^0$-复形 $Y^\bullet\in\text{obj}(\mathsf{D_0^\#\!(A)})$ 在 $\mathsf{D_0^\#\!(A)}$ 中都同构于某个 $X^\bullet\in\text{obj}(\mathsf{A})$ 的 $\mathsf{J'}(X^\bullet)$。}全体 $H^0$-复形构成的全子范畴 $\mathsf{D_0^\#\!(A)}\subset\mathsf{D^\#\!(A)}$。
\end{Pro}

\begin{proof}
	全忠实函子 $\mathsf{J}:\mathsf{A}\rightarrow\mathsf{Kom^\#\!(A)}$ 可细化为函子 $\mathsf{J}:\mathsf{A}\rightarrow\mathsf{H^\#\!(A)}$，它仍是全忠实的：$\mathsf{H^\#\!(A)}$ 中两个态射 $f^\bullet,g^\bullet$ 相等意味着 $f^\bullet\sim g^\bullet$，而由于 0-复形的微分平凡，这立刻蕴涵 $f^0=g^0$（至于在态射上的满性则是显然的）。这使得 $\mathsf{A}$ 成为 $\mathsf{H^\#\!(A)}$ 中全体 0-复形构成的全子范畴，从而对所有 $X,Y\in\text{obj}(\mathsf{A})$ 有 $\text{Hom}_\mathsf{A}(X,Y)=\text{Hom}_\mathsf{H^\#\!(A)}(X[0]^\bullet,Y[0]^\bullet)$。
	
	我们的目标是证明：对任意两个 0-复形 $X[0]^\bullet, Y[0]^\bullet$，群同态
	\[
	F:\text{Hom}_\mathsf{H^\#\!(A)}(X[0]^\bullet,Y[0]^\bullet)\rightarrow\text{Hom}_\mathsf{D^\#\!(A)}(X[0]^\bullet,Y[0]^\bullet),\,f^\bullet\mapsto [(\text{id}_{X[0]}^\bullet,f^\bullet)]
	\]
	是双射，其中 $\text{Hom}_\mathsf{D^\#\!(A)}(X[0]^\bullet,Y[0]^\bullet)=\text{Hom}_\mathsf{D^\#\!(A)}(\mathsf{L}(X[0]^\bullet),\mathsf{L}(Y[0]^\bullet))$。我们断言其逆由下式给出：
	\[
	F^{-1}:\text{Hom}_\mathsf{D^\#\!(A)}(X[0]^\bullet,Y[0]^\bullet)\rightarrow\text{Hom}_\mathsf{H^\#\!(A)}(X[0]^\bullet,Y[0]^\bullet),\; [(s^\bullet,f^\bullet)]\mapsto f^\bullet\;.
	\]
	显然，$F^{-1}\circ F$ 在 $\text{Hom}_\mathsf{H^\#\!(A)}(X[0]^\bullet,Y[0]^\bullet)$ 上是恒等。反之，给定 $\varphi=[(s^\bullet,f^\bullet)]\in\text{Hom}_\mathsf{D^\#\!(A)}(X[0]^\bullet,Y[0]^\bullet)$，令 $\phi\coloneqq (F\circ F^{-1})(\varphi)=[(\text{id}_{X[0]}^\bullet,g^\bullet)]\in\text{Hom}_\mathsf{D^\#\!(A)}(X[0]^\bullet,Y[0]^\bullet)$，其中 $g^\bullet$ 是 $\mathsf{H^\#\!(A)}$ 中的态射，由 $H^\#\!(\mathsf{A})$ 中的 $g^0\coloneqq H^0(f^\bullet)\circ H^0(s^\bullet)^{-1}:X\cong H^0(X[0]^\bullet)\rightarrow H^0(Y[0]^\bullet)\cong Y$ 给出。按屋顶图的等价关系，$\phi=\varphi$ 在 $\mathsf{D^\#\!(A)}$ 中成立当且仅当能使
	\[
	\begin{tikzcd}
		& & V^\bullet\arrow[dl, "r^\bullet"']\arrow[dr, "h^\bullet"] & & \\
		& Z^\bullet\arrow[dl, "s^\bullet"'] & & X[0]^\bullet\arrow[dlll, "{\text{id}_{X[0]}^\bullet}"]\arrow[dr, "g^\bullet"] & \\
		X[0]^\bullet & & & & Y[0]^\bullet\arrow[from=ulll, crossing over, "f^\bullet"'] 
	\end{tikzcd}
	\]
	经由 $\mathsf{D^\#\!(A)}$ 中某个适当的覆盖屋顶图而交换。这由以下数据给出：$V^{-n}\coloneqq Z^{-n}$，$V^0\coloneqq\ker(d_Z^0)$，$n>0$ 时 $V^n\coloneqq 0$（微分取诱导的），以及自然嵌入 $r^\bullet:V^\bullet\rightarrow Z^\bullet$（它是拟同构，因为在非正次数上与恒等在上同调中相同），还有除 $h^0\coloneqq H^0(s^\bullet)\circ q: \ker(d_Z^0)\rightarrow H^0(Z^\bullet)\rightarrow X$ 之外皆为零的 $h^\bullet$。交换性容易验证：对左三角形，唯一非平凡的等式是对所有 $z\in V^0$ 有 $h^0(z)=H^0(s^\bullet)\langle z\rangle=s^0(z)=(r^0\circ s^0)(z)$；对右三角形则是 $(g^0\circ h^0)(z)=H^0(f^\bullet)\langle z\rangle=f^0(z)=(f^0\circ r^0)(z)$。因此，两个屋顶图都表示 $\varphi=\phi$，这意味着 $\mathsf{F}\circ\mathsf{F}^{-1}$ 在 $\text{Hom}_\mathsf{D^\#\!(A)}(X[0]^\bullet,Y[0]^\bullet)$ 上是恒等，从而 $\mathsf{L}$ 是全忠实的。于是 $\mathsf{J}'$ 确实是全忠实的（作为全忠实函子的复合）。
	
	最后，我们证明任意 $H^0$-复形 $Z^\bullet$ 在 $\mathsf{D_0^\#\!(A)}$ 中都同构于某个 0-复形 $X[0]^\bullet$。显然的候选者是它的上同调对应物 $X[0]^\bullet\coloneqq\mathsf{J}'(H^0(Z^\bullet))=(...\rightarrow 0\rightarrow H^0(Z^\bullet)\rightarrow 0\rightarrow...)$（由 $H^0$-复形的定义，这是合理的选择）。同构正是由上面图中的覆盖屋顶图给出，其中 $h^\bullet$ 是拟同构，因为 $r^\bullet$、$s^\bullet$ 与 $\text{id}_{X[0]}^\bullet$ 都是拟同构。因此，若 $\mathsf{J}'$ 限制到其像上，便给出 $\mathsf{A}$ 与全子范畴 $\mathsf{D_0^\#\!(A)}$ 之间的等价。 
\end{proof}

% [Tutor Note: 含糊点拨] 证明中「交换性容易验证」处，作者实际只写下了 $V^0$ 上两个三角形各一个等式，其余分量因 $V^n=0$（$n>0$）或 $V^\bullet$ 与 $Z^\bullet$ 在非正次数重合而平凡；关键是由 $V^0=\ker(d_Z^0)$ 知 $d_Z^0$ 在 $V^0$ 上为零，从而 $h^\bullet$ 自动是链映射。另注意这里用到「拟同构在三项中两项为同构时第三项亦为同构」这一事实，它来自五引理。

因此，我们有时会省略函子 $\mathsf{J}'$，直接把 $\mathsf{A}$ 理解为 $\mathsf{A}\subset\mathsf{D^\#\!(A)}$。


\subsection{Ext-阿贝尔群}\label{ch2.6}


\begin{Def}
	设 $\mathsf{A}$ 为阿贝尔范畴，$X, Y\in\text{obj}(\mathsf{A})$，$k\in\mathbb{Z}$ 任意。则第 $k$ 个\textbf{Ext-阿贝尔群}\index{Ext-abelian group}由
	\begin{equation}\label{Extgroup}
		\text{Ext}_\mathsf{A}^k(X,Y)\coloneqq \text{Hom}_{\mathsf{D(A)}}(X[0]^\bullet,Y[k]^\bullet)\,,
	\end{equation}
	给出，对任意 $l\in\mathbb{Z}$ 有 $\cong\text{Hom}_\mathsf{D(A)}(X[l]^\bullet,Y[k+l]^\bullet)\in\text{obj}(\mathsf{Ab})$——这是应用定义 \ref{translfunc} 中的平移函子 $T^l:\mathsf{D(A)}\rightarrow\mathsf{D(A)}$（一个等价）之后得到的结果，故在 Hom 群上是双射。 
	
	Hom 群的复合律进而诱导出双线性的\textbf{Yoneda 积}\index{Yoneda product}\footnote{若 $X=Y=Z$，则 $\diamond$ 使向量空间 $\bigoplus_{n\in\mathbb{Z}}\text{Ext}_\mathsf{A}^n(X,X)$ 具有分次代数结构。}
	\begin{equation}
		\diamond:\text{Ext}_\mathsf{A}^k(X,Y)\times\text{Ext}_\mathsf{A}^l(Y,Z)\rightarrow\text{Ext}_\mathsf{A}^{k+l}(X,Z)
	\end{equation}
	（由上述可知，它与平移的选取无关）。由于 $\mathsf{D(A)}$ 是加性范畴，双函子 $\text{Hom}_\mathsf{D(A)}$ 实际上给出双函子 $\text{Ext}_\mathsf{A}^k:\mathsf{A}^\text{opp}\times\mathsf{A}\rightarrow\mathsf{Ab}$，于是 $\text{Ext}_\mathsf{A}^k(X,\square)$ 是共变的，而 $\text{Ext}_\mathsf{A}^k(\square,Y)$ 是反变的。
\end{Def}

% [Tutor Note: 术语提醒] 「Ext-阿贝尔群」是作者自创的称法，主流文献径直称之为 Ext 群，记作 $\text{Ext}_\mathsf{A}^k(X,Y)$；它由定义式借助 $\mathsf{D(A)}$ 中的 Hom 集定义，因此负指标时自动为零（定理 \ref{Extprop}b.），而 $k\geq 1$ 时的消没性正刻画投射/内射对象（命题 \ref{injprojExt}）。还要注意定义式中 $X[0]^\bullet$ 与 $Y[k]^\bullet$ 的平移必须成对出现，改变平移只能同时改变两者（即 $l$ 的选取），这正是括号中「与平移选取无关」的含义。

\begin{Rem}\label{explicitExtmorph}
	当 $k>0$ 时，我们可以如下在第 $k$ 个 Ext 群中具体构造出一些态射：对\textit{无环}复形
	\begin{equation}\label{Extcomplex}
		\mkern-12mu K^\bullet\!=\!\big(...\!\rightarrow 0\rightarrow K^{-k}\!\coloneqq\! Y\rightarrow K^{-k+1}\rightarrow...\rightarrow K^0\rightarrow K^1\!\coloneqq\! X\rightarrow 0\rightarrow\!...\big)\,,
	\end{equation}
	我们令 $c_{K^\bullet}\in\text{Ext}_\mathsf{A}^k(X,Y)$ 为由屋顶图
	\begin{equation}\label{Extroof}
		\begin{tikzcd}
			& \tilde{K}^\bullet\arrow[dl, "s^\bullet"']\arrow[dr, "f^\bullet"] & \\
			X[0]^\bullet & & Y[k]^\bullet
		\end{tikzcd}\quad,
	\end{equation}
	所表示的态射，其中 $\tilde{K}^1\coloneqq 0$，其余情形 $\tilde{K}^n\coloneqq K^n$，$s^0\coloneqq d_K^0:K^0\rightarrow K^1=X$（$s^\bullet$ 中唯一的非平凡项），$f^{-k}\coloneqq\text{id}_Y:K^{-k}=Y\rightarrow Y$（$f^\bullet$ 中唯一的非平凡项）。此外，给定另一个 $l>0$ 的无环复形
	\[
	L^\bullet=\big(...\rightarrow 0\rightarrow L^{-l}\coloneqq Z\rightarrow L^{-l+1}\rightarrow...\rightarrow L^0\rightarrow L^1\coloneqq Y\rightarrow 0\rightarrow...\big)
	\]
	我们可以把 $L^\bullet$ 与 $K^\bullet$ 链接起来，得到无环复形
	\begin{small}\begin{equation}\label{astcomplex}
			L^\bullet\ast K^\bullet=\big(...\rightarrow 0\rightarrow \mkern-16mu \underbrace{L^{-l}\!=\!Z}_{=(L^\bullet\ast K^\bullet)^{-k-l}}\mkern-16mu\rightarrow \!...\!\rightarrow L^0\!\xrightarrow[d_K^{-k}\circ d_L^0]{}\! K^{-k+1}\!\rightarrow\! ...\!\rightarrow \mkern-12mu\underbrace{K^1\!=\!X}_{=(L^\bullet\ast K^\bullet)^1}\mkern-12mu\rightarrow 0\rightarrow...\big)\,,
	\end{equation}\end{small}
	\!\!它在 $\text{Ext}_\mathsf{A}^{k+l}(X,Z)$ 中对应的态射，在下面定理的最后一条中有更好的描述。
\end{Rem}

% [Tutor Note: 易错预警] 复形 \eqref{Extcomplex} 的编号是「倒着」写的：$K^{-k}=Y$ 在最低次、$K^1=X$ 在最高次，故准同构 $s^\bullet$ 只在 $0$ 次取 $d_K^0:K^0\to K^1=X$，而真正的态射 $f^\bullet$ 只在 $-k$ 次取 $\text{id}_Y$；两者次数之差恰好由目标复形的平移 $[k]$ 补上。若把复形自身的编号与 $\mathsf{D(A)}$ 中的上同调次数混为一谈，就会算错 $\text{Ext}$ 的次数。

\begin{Thm}\label{Extprop}
	设 $\mathsf{A}$ 为阿贝尔范畴，$X,Y\in\textup{obj}(\mathsf{A})$ 固定。\textup{Ext} 群具有如下性质：
	\renewcommand{\theenumi}{\alph{enumi}}
	\begin{enumerate}[leftmargin=0.5cm]
		\item $\mathsf{A}$ 中的短正合序列 $0\rightarrow X\rightarrow X'\rightarrow X''\rightarrow 0$ 与 $0\rightarrow Y\rightarrow Y'\rightarrow Y''\rightarrow 0$ 分别诱导出长正合序列
		\[
		...\rightarrow \textup{Ext}_\mathsf{A}^k(X'',Y)\rightarrow \textup{Ext}_\mathsf{A}^k(X',Y)\rightarrow \textup{Ext}_\mathsf{A}^k(X,Y)\rightarrow \textup{Ext}_\mathsf{A}^{k+1}(X'',Y)\rightarrow...\;,
		\]
		\[
		...\rightarrow \textup{Ext}_\mathsf{A}^k(X,Y)\rightarrow \textup{Ext}_\mathsf{A}^k(X,Y')\rightarrow \textup{Ext}_\mathsf{A}^k(X,Y'')\rightarrow \textup{Ext}_\mathsf{A}^{k+1}(X,Y)\rightarrow...\;.
		\]
		
		\item 对 $k<0$，$\textup{Ext}_\mathsf{A}^k(X,Y)=\{0\}$。
		
		\item 当 $k=0$ 时，$\textup{Ext}_\mathsf{A}^0(X,Y)\cong \textup{Hom}_\mathsf{A}(X,Y)$。
		
		\item 对 $k>0$ 及任意 $f\in\textup{Ext}_\mathsf{A}^k(X,Y)$，存在形如 \eqref{Extcomplex} 的某个无环复形 $K^\bullet$ 使得 $f=c_{K^\bullet}$。再给定 $c_{L^\bullet}\in\textup{Ext}_\mathsf{A}^l(Y,Z)$，则成立 $c_{L^\bullet\ast K^\bullet}=c_{L^\bullet}\diamond c_{K^\bullet}$。
	\end{enumerate}
\end{Thm}

\begin{proof}
	\ \vspace*{0.0cm}
	\renewcommand{\theenumi}{\alph{enumi}}
	\begin{enumerate}[leftmargin=0.5cm]
		\item 把所给的短正合序列看作 $\mathsf{D^b}(\mathsf{A})$ 中的区分三角形 $X^\bullet\rightarrow X'^\bullet\rightarrow X''^\bullet\rightarrow X[1]^\bullet$ 与 $Y^\bullet\rightarrow Y'^\bullet\rightarrow Y''^\bullet\rightarrow Y[1]^\bullet$（参看命题 \ref{shortexactaretriangles}）。于是我们需要的是定理 \ref{longexactcohomtriangles} 的 Hom 版本。其证明需要一些关于三角范畴（triangulated category）的知识，而导出范畴 $\mathsf{D^\#\!(A)}$ 正是其显例（见推论 \ref{D(A)istriangulated}）。我们会在适当的时候回到这一点……
		
		\item 对 $k>0$，考虑由实线屋顶图
		\begin{equation}\label{commutingroofs}
			\begin{tikzcd}
				& X'^\bullet\arrow[dl, "s^\bullet"']\arrow[dr, "f^\bullet"] & \\
				X[0]^\bullet & & Y[-k]^\bullet \\
				& \tilde{X}^\bullet\arrow[ul, dashed, "t^\bullet"]\arrow[uu, dashed, "r^\bullet"]\arrow[ur, dashed, "0"'] 
			\end{tikzcd}\quad.
		\end{equation}
		表示的态射 $\varphi=[(s^\bullet,f^\bullet)]\in\text{Ext}_\mathsf{A}^{-k}(X^\bullet,Y^\bullet)=$ $\text{Hom}_\mathsf{D(A)}(X[0]^\bullet,Y[-k]^\bullet)$。用 $\tilde{X}^j\coloneqq X'^j$（当 $j<k-1$）、$\tilde{X}^{k-1}\coloneqq\ker(d_{X'}^{k-1})$、其余情形 $\tilde{X}^j\coloneqq 0$ 定义新的子复形 $\tilde{X}^\bullet\xhookrightarrow{r^\bullet}X'^\bullet$，其微分由 $X'^\bullet$ 上的微分诱导。再令 $t^0\coloneqq s^0$，作为 $t^\bullet$ 的唯一非平凡项。这些数据纳入图中的虚线部分，而该图随即显然交换（特别地，$s^0\circ r^0=s^0=t^0$，且由于 $\tilde{X}^k=0$ 而有 $f^k\circ r^k=0$）。现在，由于按定义 $H^n(s^\bullet)$ 总是同构，我们有 $H^n(X'^\bullet)=0$，除非 $H^0(X'^\bullet)\cong X$。这一构造必然蕴涵 $r^\bullet$ 与 $t^\bullet$ 都是拟同构。我们于是断定下方虚线屋顶图是 $\varphi$ 的一个有效代表元……但它正是零屋顶图，故 $\varphi=0$，再由其任意性得 $\text{Ext}_\mathsf{A}^{-k}(X^\bullet,Y^\bullet)=\{0\}$，如所愿。
		
		\item 现在设 $k=0$。在命题 \ref{AinD(A)} 的证明中，我们已经证明函子 $\mathsf{J}'\!:\!\mathsf{A}\rightarrow\mathsf{D(A)}$ 是全忠实的，故 $\text{Hom}_\mathsf{A}(X,Y)\cong\text{Hom}_\mathsf{D(A)}(X[0]^\bullet,Y[0]^\bullet)$ $=\text{Ext}_\mathsf{A}^0(X^\bullet,Y^\bullet)$。
		
		\item 关于 $k>0$ 的陈述，其证明更为繁复：大致做法是从一个初始屋顶图出发逐次细化，使其越来越接近所期望的图 \eqref{Extroof}，而每一步的相容性由形如 \eqref{commutingroofs} 的交换图保证。有兴趣的读者可参看 \cite[Theorem III.5.5]{[GM03]} 的证明的 c) 部分。 
		\newline 至于与乘积相容的行为，设 $c_{K^\bullet}=[(s^\bullet,f^\bullet)]\in\text{Ext}_\mathsf{A}^k(X,Y)$、$c_{L^\bullet}=[(r^\bullet,g^\bullet)]\in\text{Ext}_\mathsf{A}^l(Y,Z)$ 与 $c_{L^\bullet\ast K^\bullet}=[(t^\bullet,h^\bullet)]\in\text{Ext}_\mathsf{A}^{k+l}(X,Z)$ 都如 \eqref{Extroof} 那样由屋顶图表示。考虑屋顶图的（Yoneda）乘积法则：
		\[
		\begin{tikzcd}[column sep=0.5cm]
			& & \widetilde{L^\bullet\ast K^\bullet}\arrow[dl, dashed, "\bar{t}^\bullet"]\arrow[dr, dashed, "\bar{h}^\bullet"']\arrow[ddll, out=195, in=68,  "t^\bullet"']\arrow[ddrr, out=-15, in = 120, "h^\bullet"] & & \\
			& \tilde{K}^\bullet\arrow[dl, "s^\bullet"]\arrow[dr, "f^\bullet"'] & & \tilde{L}[k]^\bullet\arrow[dl, "{r[k]^\bullet}"]\arrow[dr, "{g[k]^\bullet}"'] & \\
			X[0]^\bullet & & Y[k]^\bullet & & Z[k+l]^\bullet
		\end{tikzcd}\quad.
		\]
		参照链接起来的复形 \eqref{astcomplex}，用 $\bar{h}^i\coloneqq\text{id}_{L^{i+1}}$（当 $-k-l\leq i\leq -k$）定义虚线箭头，相应地 $\bar{t}^i\coloneqq\text{id}_{K}^i$（当 $-k+1\leq i\leq 0$）以及 $\bar{t}^0\coloneqq d_L^0$（其余为零）。确实有 $s^\bullet\circ\bar{t}^\bullet=t^\bullet$、$g[k]^\bullet\circ\bar{h}^\bullet= h^\bullet$ 与 $f^\bullet\circ\bar{t}^\bullet=r[k]^\bullet\circ\bar{h}^\bullet$，这证明了 $c_{L^\bullet}\diamond c_{K^\bullet}=c_{L^\bullet\ast K^\bullet}$。   
	\end{enumerate}\vspace*{-0.5cm}
\end{proof}

% [Tutor Note: 记号提醒] 证明 b. 条时作者写「对 $k>0$」，但所考察的对象是 $\text{Ext}_\mathsf{A}^{-k}$，即负指标；定理 b. 的陈述本身是「$k<0$ 时为零」，两处相差一个负号，阅读时请以实际指标 $-k$ 为准。另外此处把变量写成 $X^\bullet,Y^\bullet$（带上标），而定义式用的是 $X[0]^\bullet,Y[k]^\bullet$，二者在上下文中视为同一对象。

下面给出几个定义，用以刻画导出范畴中态射的丰富性。它们是为分析提供的一个额外工具，但在本文其余部分实际上不会起作用。

\begin{Def}
	设 $\mathsf{A}$ 为阿贝尔范畴。我们定义其\textbf{同调维数}（homological dimension）\index{homological dimension}为\footnote{注意到在「最坏」情形下 $\text{dim}_h(\mathsf{A})=0$，因为由定理 \ref{Extprop}c.，对任意 $X\in\text{obj}(\mathsf{A})$ 有 $\text{Ext}_\mathsf{A}^0(X,X)=\text{Hom}_\mathsf{A}(X,X)\ni\text{id}_X$。对 $\text{dim}_p(X)$、$\text{dim}_i(X)$ 类似。} 
	\begin{equation}\label{homologicaldimension}
		\text{dim}_h(\mathsf{A})\coloneqq\text{max}\big(\{k\in\mathbb{N}_0\mid \exists X,Y\in\text{obj}(\mathsf{A})\text{ s.t. }\text{Ext}_\mathsf{A}^k(X,Y)\neq\{0\}\}\big)
	\end{equation}
	（若不存在这样的 $k$，则 $=\infty$）。类似地，固定对象 $X\in\text{obj}(\mathsf{A})$ 的\textbf{投射/内射同调维数}\index{homological dimension!projective/injective}为 
	\begin{equation}
		\begin{aligned}
			\text{dim}_p(X)&\coloneqq\text{sup}\big(\{k\in\mathbb{N}_0\mid \exists Y\in\text{obj}(\mathsf{A})\text{ s.t. }\text{Ext}_\mathsf{A}^k(X,Y)\neq\{0\}\}\big) \\
			\text{dim}_i(X)&\coloneqq\text{sup}\big(\{k\in\mathbb{N}_0\mid \exists Y\in\text{obj}(\mathsf{A})\text{ s.t. }\text{Ext}_\mathsf{A}^k(Y,X)\neq\{0\}\}\big)
		\end{aligned}
	\end{equation}
	（同样可能为 $\infty$）。    
\end{Def}

% [Tutor Note: 含糊点拨] 脚注中「最坏情形 $\text{dim}_h(\mathsf{A})=0$」实指维数取到最小可能值 $0$：因 $\text{Ext}_\mathsf{A}^0(X,X)=\text{Hom}_\mathsf{A}(X,X)\ni\text{id}_X$ 恒非零，定义式中取 max 的集合总是含 $0$，故「不存在这样的 $k$」在定义本身的逻辑下不会发生，括号中的 $=\infty$ 只是形式上的约定；真正「最坏」（即维数最大、Ext 无消没）的情形才对应无限整体维数。

\begin{Pro}\label{injprojExt}
	设 $\mathsf{A}$ 为阿贝尔范畴。以下关于同调维数的条件是等价的：
	\renewcommand{\theenumi}{\roman{enumi}}
	\begin{enumerate}[leftmargin=0.5cm]
		\item $\textup{dim}_h(\mathsf{A})=0 \xLeftrightarrow{(1)} \textup{Ext}_\mathsf{A}^1(X,Y)\!=\!\{0\}$ $\forall X,Y\!\in\!\textup{obj}(\mathsf{A}) \xLeftrightarrow{(2)} \mathsf{A}$ 是半单的。
		
		\item $\textup{dim}_p(X)=0 \xLeftrightarrow{(1)} \textup{Ext}_\mathsf{A}^1(X,Y)=\{0\}$ $\forall Y\!\in\!\textup{obj}(\mathsf{A}) \xLeftrightarrow{(2)} X$ 是投射对象（projective object）。
		
		\item $\textup{dim}_i(X)=0 \xLeftrightarrow{(1)} \textup{Ext}_\mathsf{A}^1(Y,X)=\{0\}$ $\forall Y\!\in\!\textup{obj}(\mathsf{A}) \xLeftrightarrow{(2)} X$ 是内射对象（injective object）。
	\end{enumerate}
\end{Pro}

\begin{proof}
	\ \vspace*{0.0cm}
	\renewcommand{\theenumi}{\roman{enumi}}
	\begin{enumerate}[leftmargin=0.5cm]
		\item 蕴涵 $\xRightarrow{(1)}$ 是显然的。至于 $\xLeftarrow{(1)}$，考虑 $k>1$ 及按注 \ref{explicitExtmorph} 构造的任意 $f=c_{K^\bullet}\in\text{Ext}_\mathsf{A}^k(X,Y)$。则两个复形
		\begin{align*}
			\mkern-12mu K_1^\bullet\!&\coloneqq\!\big(...\!\rightarrow 0\rightarrow K^{-k}\!\coloneqq\! Y\rightarrow...\rightarrow K^{-1}\rightarrow Z\!\coloneqq\!\text{im}(d_K^{-1})\rightarrow 0\rightarrow\!...\big)\,, \\
			K_2^\bullet\!&\coloneqq\!\big(...\!\rightarrow 0\rightarrow \text{im}(d_K^{-1})\rightarrow K^0\rightarrow K^1\coloneqq X\rightarrow 0\rightarrow ...\big)
		\end{align*}
		给出 $K_1^\bullet\ast K_2^\bullet=K^\bullet$，故由定理 \ref{Extprop}d. 有 $f=c_{K_1^\bullet}\diamond c_{K_2^\bullet}$。但 $c_{K_2^\bullet}$ $\in\text{Ext}_\mathsf{A}^1(X,Z)=\{0\}$，故 $f=0$，从而 $\text{Ext}_\mathsf{A}^k(X,Y)=\{0\}$。
		\newline 至于 $\xRightarrow{(2)}$，考虑短正合序列 $0\rightarrow Y\xrightarrow{f} Z\xrightarrow{g} X\rightarrow 0$，由定理 \ref{Extprop}a.\&b. 它产生（长）正合序列的一部分 
		\vspace*{-0.2cm}
		\[
		\{0\}\rightarrow\text{Hom}_\mathsf{A}(X,Y)\xrightarrow{f_*}\text{Hom}_\mathsf{A}(X,Z)\xrightarrow{g_*}\text{Hom}_\mathsf{A}(X,X)\rightarrow\text{Ext}_\mathsf{A}^1(X,Y)\!=\!\{0\}.
		\]
		于是存在某个 $h:X\rightarrow Z$ 使 $g_*(h)=g\circ h=\text{id}_X$，它还可经由典范同构 $k:X\rightarrow\text{im}(h)$ 与包含 $\iota:\text{im}(h)\rightarrow Z$ 分解。五引理的直接应用告诉我们，$0\rightarrow Y\xrightarrow{(\text{id}_Y,0)} Y\oplus\text{im}(h)\xrightarrow{0+k^{-1}} X\cong\text{im}(h)\rightarrow 0$ 同构于原来的序列（其中间同构为 $f+\iota:Y\oplus\text{im}(h)\rightarrow Z$，它使相应的图交换）。因此，任意短正合序列都分裂，这正是 $\mathsf{A}$ 半单的定义（参看注 \ref{semisimple}）。反之，$\text{Ext}_\mathsf{A}^1(X,Y)$ 中的任意态射都由三元组 $K^\bullet=(0\rightarrow Y\xrightarrow{f} K^0\xrightarrow{g} X\rightarrow 0)$ 构成，依假设它分裂，最终可证明这给出 $c_{K^\bullet}=0$，于是得蕴涵 $\xLeftarrow{(2)}$。
		
		\item $\xRightarrow{(1)}$ 是清楚的；$\xLeftarrow{(1)}$ 如 i. 中那样，通过分裂式 $\text{Ext}_\mathsf{A}^k(X,Y)=\text{Ext}_\mathsf{A}^{k-1}(\text{im}(d_K^{-1}),Y)\diamond\text{Ext}_\mathsf{A}^1(X,\text{im}(d_K^{-1}))=0$（对所有 $k\geq 1$）来证明。
		\newline 关于 $\xLeftrightarrow{(2)}$，i. 中 $\xLeftrightarrow{(2)}$ 的相应论断蕴涵：给定某个固定的 $X\in\text{obj}(\mathsf{A})$，对所有 $Y\in\text{obj}(\mathsf{A})$ 有 $\text{Ext}_\mathsf{A}^1(X,Y)=\{0\}$，当且仅当任何以 $X$ 为第三个对象的短正合序列都分裂。另外回忆 $X$ 的投射性条件，即图 \eqref{injprojobjdiag}，\vspace*{-0.5cm} 
		\[
		\begin{tikzcd}
			& & & X\arrow[dl, dashed, "h"']\arrow[d, "f"] & \\
			\Big(0\arrow[r] & Y\arrow[r, hook] & \Big)\;Z\arrow[r, two heads, "g"'] & U\arrow[r] & 0
		\end{tikzcd}\vspace*{-0.3cm}
		\]
		（实线箭头为已知数据）。至于 $\xRightarrow{(2)}$，令 $Z\times_U X\!\coloneqq\!\{(z,x)\in Z\!\times\! X\mid g(z)\!=\!f(x)\}$。由投影 $p_2:Z\times_U X\rightarrow X$ 的满射性以及所出现的短正合序列依假设分裂，存在包含 $i_2:X\rightarrow Z\times_U X$。于是 $h\coloneqq p_1\circ i_2:X\rightarrow Z$ 使图交换，证明 $X$ 是投射对象。反之对 $\xLeftarrow{(2)}$，在上式中取 $U=X$ 与 $f=\text{id}_X$，则依假设存在 $h:X\rightarrow Z$ 使 $g\circ h=\text{id}_X$，这说明任意短正合序列 $0\rightarrow Y\hookrightarrow Z \twoheadrightarrow X\rightarrow 0$ 都同构于以 $Y\oplus\text{im}(h)$ 为中间项的分裂序列。如上所述，这等价于对所有 $Y\in\text{obj}(\mathsf{A})$ 有 $\text{Ext}_\mathsf{A}^1(X,Y)=\{0\}$。
		
		\item 这一条与前一部分类似地证明。 
	\end{enumerate}\vspace*{-0.5cm}
\end{proof}

下面是另一个涉及投射对象与内射对象的有用结果。

\begin{Lem}
	设 $\mathsf{A}$ 为阿贝尔范畴，$X,X'\in\textup{obj}(\mathsf{A})$。
	\renewcommand{\theenumi}{\roman{enumi}}
	\begin{enumerate}[leftmargin=0.5cm]
		\item 给定无环复形
		\[
		P^\bullet\coloneqq\big(...\rightarrow 0\rightarrow X'\rightarrow P^{-k}\rightarrow...\rightarrow P^0\rightarrow X\rightarrow 0\rightarrow...\big)\,,
		\]
		其中 $P^i\in\textup{obj}(\mathsf{A})$ 是投射对象，则成立 $\textup{dim}_p(X')=\textup{max}(\{\textup{dim}_p(X)-k-1,0\})$。
		
		\item 给定无环复形
		\[
		I^\bullet\coloneqq\big(...\rightarrow 0\rightarrow X\rightarrow I^0\rightarrow...\rightarrow I^k\rightarrow X'\rightarrow 0\rightarrow...\big)\,,
		\]
		其中 $I^i\in\textup{obj}(\mathsf{A})$ 是内射对象，则成立 $\textup{dim}_i(X')=\textup{max}(\{\textup{dim}_i(X)-k-1,0\})$。 
	\end{enumerate}
\end{Lem}

\begin{proof}
	我们证明投射情形 i. 的陈述（内射情形可用类似的步骤证明）。给定任意 $Y\in\text{obj}(\mathsf{A})$，我们通过对 $k$ 的归纳构造映射 $\delta:\text{Ext}_\mathsf{A}^d(X',Y)\rightarrow\text{Ext}_\mathsf{A}^{d+k+1}(X,Y)$，它在 $d=0$ 时是满射，在 $d\geq 1$ 时是同构。
	
	当 $k=0$ 时，$P^\bullet=(0\rightarrow X'\rightarrow P^0\rightarrow X\rightarrow 0)$ 由定理 \ref{Extprop}a. 诱导出长正合序列的一部分
	\[
	\text{Ext}_\mathsf{A}^d(P^0,Y)\rightarrow\text{Ext}_\mathsf{A}^d(X',Y)\xrightarrow{\delta}\mkern-15mu\rightarrow\text{Ext}_\mathsf{A}^{d+1}(X,Y)\rightarrow\text{Ext}_\mathsf{A}^{d+1}(P^0,Y)=\{0\}\,,
	\]
	于是 $\delta$ 显然是满射，并且当 $d>0$ 时还是同构（因为由命题 \ref{injprojExt}ii.，$P^0$ 为投射对象故 $\text{Ext}_\mathsf{A}^d(P^0,Y)=\{0\}$；第三项同理）。至于归纳步 $k-1\rightarrow k$，把 $P^\bullet=P_1^\bullet\ast P_2^\bullet$ 拆成 $P_1^\bullet\coloneqq\big(0\rightarrow X'\rightarrow P^{-k}\rightarrow X''\coloneqq\text{im}(d_P^{-k})\rightarrow 0\big)$ 与 $P_2^\bullet\coloneqq\big(0\rightarrow X''\rightarrow\break P^{-k+1}\rightarrow...\rightarrow P^0\rightarrow X\rightarrow 0\big)$。对 $P_2^\bullet$ 的归纳假设给出 $\delta_2:\text{Ext}_\mathsf{A}^d(X'',Y)\rightarrow\text{Ext}_\mathsf{A}^{d+k}(X,Y)$，而 $\delta_1:\text{Ext}_\mathsf{A}^d(X',Y)\rightarrow\text{Ext}_\mathsf{A}^{d+1}(X'',Y)$ 如同基始情形那样由 $P_1^\bullet$ 推出。则 $\delta\coloneqq\delta_2[1]\circ\delta_1$ 继承所期望的性质。
	
	我们断定：若 $d\coloneqq\text{dim}_p(X')\geq 1$，即存在某个 $Y\in\text{obj}(\mathsf{A})$ 使 $\text{Ext}_\mathsf{A}^{\text{dim}_p(X')+k+1}(X,Y)\cong\text{Ext}_\mathsf{A}^{\text{dim}_p(X')}(X',Y)\neq\{0\}$，则 $\text{dim}_p(X)=\text{dim}_p(X')+k+1$；而 $\text{dim}_p(X')= 0$ 蕴涵对所有 $Y$ 有 $\text{Ext}_\mathsf{A}^1(X',Y)=\{0\}$，于是 $\text{Ext}_\mathsf{A}^{k+2}(X,Y)=\{0\}$ 且 $\text{dim}_p(X)\leq k+1$。合起来即得 $\text{dim}_p(X')=\text{max}(\{\text{dim}_p(X)-k-1,0\})$。 
\end{proof}

\begin{Cor}
	设 $\mathsf{A}$ 为阿贝尔范畴，$X\in\textup{obj}(\mathsf{A})$。
	\begin{itemize}[leftmargin=0.5cm]
		\item 若 $\mathsf{A}$ 有足够多的投射对象\textup(见定义 \ref{injprojobj}\textup)，则 $\textup{dim}_p(X)\leq k$ 当且仅当存在 $X$ 的长度 $\leq k+1$ 的投射消解（projective resolution）。
		
		\item 若 $\mathsf{A}$ 有足够多的内射对象，则 $\textup{dim}_i(X)\leq k$ 当且仅当存在 $X$ 的长度为 $k+1$ 的内射消解（injective resolution）。 
	\end{itemize} 
\end{Cor}

% [Tutor Note: 易错预警] 两条陈述中的长度界不对称：第一条为「长度 $\leq k+1$」，第二条为「长度为 $k+1$」。按通行的约定（消解的长度从 $0$ 次起算，$\text{dim}_p(X)\leq k$ 对应长度 $\leq k+1$ 的投射消解），两条应作同样理解，引用时请以第一条的 $\leq k+1$ 为准。

我们还要提到一个引理，它对我们第 7 章中椭圆曲线的分析至关重要。它出自 \cite[Exercise 13.13]{[KS06]}。

\begin{Lem}\label{whendimhleq1}
	设 $\mathsf{A}$ 为满足 $\textup{dim}_h(\mathsf{A})\leq 1$ 的阿贝尔范畴。则任意 $X^\bullet\in\textup{obj}(\mathsf{D^b(A)})$ 都可写成 $X^\bullet\cong\bigoplus_{n\in\mathbb{Z}} H^n(X^\bullet)[-n]^\bullet\in\textup{obj}(\mathsf{D^b(A)})$。
\end{Lem}


\newpage



